% !TEX program = xelatex
\documentclass[UTF8,10pt,a4paper,fontset=none]{ctexart}

\usepackage[a4paper,left=19mm,right=19mm,top=18mm,bottom=19mm,headheight=14pt,headsep=7mm,footskip=10mm]{geometry}
\usepackage{fontspec}
\usepackage{amsmath}
\usepackage{booktabs,longtable,tabularx,array}
\usepackage{enumitem}
\usepackage{xcolor}
\usepackage{graphicx}
\usepackage{caption}
\usepackage{hyperref}

\setmainfont{Times New Roman}
\setsansfont{Arial}
\setmonofont{Menlo}[Scale=0.86]
\setCJKmainfont[Script=Default]{Songti SC}
\setCJKsansfont[Script=Default]{Songti SC}
\setCJKmonofont[Script=Default]{Songti SC}
\DeclareMathSizes{10.95}{10}{10}{10}
\DeclareMathSizes{10}{10}{10}{10}
\DeclareMathSizes{9}{10}{10}{10}
\DeclareMathSizes{8}{10}{10}{10}
\DeclareMathSizes{7}{10}{10}{10}

\definecolor{deepblue}{HTML}{17365D}
\definecolor{teal}{HTML}{2459A6}
\definecolor{softblue}{HTML}{EAF1FB}
\definecolor{softgray}{HTML}{F4F6F8}
\definecolor{rulegray}{HTML}{D7DFEA}
\definecolor{codegreen}{HTML}{2E6B4F}

\hypersetup{
  colorlinks=true,
  linkcolor=deepblue,
  citecolor=teal,
  urlcolor=deepblue,
  pdftitle={LLM Scaling Law：从经验幂律、Compute-Optimal 到数据与推理扩展},
  pdfauthor={Jiguo Li},
  pdfsubject={LLM Scaling Law 技术报告},
  pdfkeywords={LLM, Scaling Law, Compute-Optimal, Data Scaling, MoE, Inference-Time Compute}
}

\linespread{1.13}
\setlength{\parindent}{2em}
\setlength{\parskip}{2.5pt}
\setlength{\abovedisplayskip}{6pt}
\setlength{\belowdisplayskip}{6pt}
\setlength{\abovedisplayshortskip}{4pt}
\setlength{\belowdisplayshortskip}{4pt}
\setlength{\jot}{3pt}
\emergencystretch=2em
\setlist{nosep,leftmargin=2em,topsep=2pt,itemsep=1pt,parsep=0pt}
\renewcommand{\arraystretch}{1.16}
\captionsetup{font=small,labelfont=bf,labelsep=quad,hypcap=false}
\ctexset{
  section={format=\zihao{3}\bfseries\sffamily\color{deepblue},beforeskip=1.8ex,afterskip=0.65ex},
  subsection={format=\zihao{4}\bfseries\sffamily\color{deepblue},beforeskip=1.35ex,afterskip=0.45ex},
  subsubsection={format=\normalsize\bfseries\sffamily\color{deepblue},beforeskip=1.0ex,afterskip=0.3ex}
}

\newcommand{\reportshorttitle}{LLM Scaling Law 技术报告}
\makeatletter
\def\ps@scalingreport{%
  \def\@oddhead{\vbox{\hbox to\textwidth{\small\color{gray}\reportshorttitle\hfill\leftmark}\vskip3pt\hrule height0.3pt}}%
  \let\@evenhead\@oddhead
  \def\@oddfoot{\hfil\small\thepage\hfil}%
  \let\@evenfoot\@oddfoot
}
\makeatother
\pagestyle{scalingreport}
\renewcommand{\sectionmark}[1]{\markboth{\thesection\quad #1}{}}

\newcolumntype{Y}{>{\raggedright\arraybackslash}X}
\newcolumntype{C}[1]{>{\centering\arraybackslash}p{#1}}
\newcolumntype{R}[1]{>{\raggedleft\arraybackslash}p{#1}}
\newcommand{\codeid}[1]{\textnormal{#1}}
\newcommand{\evidence}[1]{\textcolor{teal}{\textbf{#1}}}
\newcommand{\term}[1]{\textbf{#1}}
\newsavebox{\fitcodebox}

\newcommand{\noteBox}[2]{%
  \begin{center}
  \setlength{\fboxsep}{7pt}%
  \setlength{\fboxrule}{0.4pt}%
  \fcolorbox{teal}{softblue}{\parbox{0.92\linewidth}{%
    \textbf{\color{deepblue}#1：}\color{black}#2}}
  \end{center}}

\begin{document}

\begin{center}
  {\sffamily\bfseries\zihao{1}\color{deepblue}
  LLM Scaling Law：从经验幂律、\par Compute-Optimal 到数据与推理扩展}

  \vspace{0.6em}
  {\zihao{3}\sffamily 一份面向预训练实践的技术认知与判断指南}

  \vspace{0.55em}
  {\normalsize\textbf{Jiguo Li\footnote{本文在 Codex 协助下完成}}}\par

  \vspace{0.2em}
  {\small\href{mailto:jiguolee@gmail.com}{jiguolee@gmail.com}}\par

  \vspace{0.3em}
  {\small\sffamily 调研与修订日期：2026 年 8 月 7 日}
\end{center}

\vspace{0.25em}
\noindent\color{deepblue}\rule{\textwidth}{0.8pt}\color{black}
\vspace{0.3em}
\begin{abstract}
LLM scaling law 不是“参数越大、能力越强”的口号，而是在固定数据分布、模型族与训练配方下，用一组受控的小规模实验刻画损失随参数量、训练 token 与计算量变化的局部经验规律，并据此进行资源配置和外推决策。本文沿着 learning curve、Kaplan 三轴幂律、Chinchilla compute-optimal、data-constrained/data-quality scaling、MoE 与 inference-time scaling 的脉络，给出核心公式、推导、拟合和验证细节。报告进一步整理 \textbf{2020--2026 年 28 个公开模型家族}的尺寸、训练 token、total/active parameters 与证据等级，用来检验公开实践支持什么、不能支持什么。面向预训练工程，本文给出 isoFLOP 实验矩阵、鲁棒拟合、scale holdout、bootstrap、不确定性传播，以及 code/repo-level 数据策略的可量化验证框架。核心结论是：\textbf{同分布 validation loss 的局部外推最可靠}；\textit{跨代 benchmark 与能力涌现最不可靠}；\textbf{函数形式可以迁移，系数、最优 token/parameter 比和能力阈值通常不能直接迁移}。

\vspace{0.5em}
\noindent\textbf{关键词：}大语言模型；Scaling Law；Compute-Optimal；数据扩展；MoE；推理时计算；代码预训练
\end{abstract}

\noteBox{核心判断}{Scaling law 是\term{局部经验模型、受控实验方法与资源决策工具}的组合。它最擅长回答“在相同 recipe 下，扩大多少参数或 token 能降低多少 loss”，而不是回答“某种智能何时必然出现”。}

\clearpage
\pdfbookmark[0]{目录}{toc}
\begingroup
\setcounter{tocdepth}{2}
\small
\setlength{\parskip}{0pt}
\linespread{1.02}\selectfont
\tableofcontents
\endgroup
\clearpage

\section{引言：先分清四类 Scaling 问题}

大模型讨论中，至少有四类问题经常被混在一起。第一类是\term{loss scaling}：模型参数量 $N$、训练 token 数 $D$ 或训练计算量 $C$ 增长时，held-out cross-entropy loss 是否在有限尺度区间内平滑下降。第二类是\term{compute-optimal scaling}：给定训练预算，应把计算分给更大的模型还是更多的数据。第三类是\term{capability scaling}：连续的 loss 改善是否会变成连续的 MMLU、代码生成或推理能力。第四类是\term{system/economic scaling}：训练 FLOP-optimal 是否也是推理和部署成本最优。

这四个问题的证据强度不同。连续、同分布且与训练目标一致的 validation loss 通常噪声最低；exact match、pass@1 等离散指标受到阈值、prompt、采样和污染影响；全生命周期最优又取决于调用量、显存、延迟和服务生命周期。因此，本文采用表~\ref{tab:evidence-level} 的证据分级。

\begin{table}[htbp]
\centering
\caption{Scaling 证据分级。只有 A 级证据适合拟合规律，B/C 级主要用于趋势校验。}
\label{tab:evidence-level}
\begin{tabularx}{\textwidth}{C{1.1cm}YY}
\toprule
等级 & 定义 & 典型材料 \\
\midrule
\evidence{A} & 同架构、同 tokenizer、同数据分布，系统改变 $N/D/C$，公开 loss 或 checkpoints & Chinchilla pilot、Pythia、Cerebras-GPT \\
\evidence{B} & 同代模型家族，关键训练信息公开，但尺寸点较少或 token/recipe 随尺寸改变 & GPT-3、LLaMA、Llama 2/3.1、Qwen2.5 \\
\evidence{C} & 跨代或单 endpoint，架构、数据、多模态或 post-training 改动显著 & Gemma 2/4、Llama 4、DeepSeek-V4 \\
\bottomrule
\end{tabularx}
\end{table}

\noteBox{关键证据边界}{公开模型尺寸表不是一条 universal scaling curve。不同 tokenizer、数据质量、代码比例、架构、上下文长度和 post-training 都是强混杂变量；跨家族散点最多说明工程趋势，不能识别参数量的因果效应。}

\section{问题起源：从 Learning Curve 到语言模型幂律}

\subsection{早期起点：误差随数据量呈幂律下降}

机器学习长期使用 learning curve 观察数据量与泛化误差的关系。现代神经 scaling law 的转折点是 Hestness 等人在机器翻译、语言建模、图像和语音等任务上观察到：在多个数量级内，可约误差常近似服从简单幂律，最优模型规模也随数据量规律增长\cite{hestness2017}。Rosenfeld 等人进一步提出同时依赖模型规模和数据规模的函数形式，并用较小尺度实验预测更大尺度表现\cite{rosenfeld2019}。

单变量幂律通常写为
\begin{equation}
  L(x)=L_{\infty}+A x^{-\alpha},
  \label{eq:single-power}
\end{equation}
其中 $x$ 可以是参数量、数据量或计算量；$L_{\infty}$ 是当前分布与模型族下的渐近损失下限；$A x^{-\alpha}$ 是可通过扩展资源降低的部分；$\alpha>0$ 是扩展效率。减去 $L_{\infty}$ 后，式~\eqref{eq:single-power} 在 log--log 坐标中成为直线：
\begin{equation}
  \log\!\left(L-L_{\infty}\right)=\log A-\alpha\log x.
  \label{eq:loglinear}
\end{equation}
指数通常很小，意味着边际收益递减很强。更重要的是，$L_{\infty}$ 未知时不能直接做普通线性回归；否则斜率会产生系统偏差。

\subsection{Kaplan：统一参数、数据和计算三轴}

Kaplan 等人把语言模型的参数量 $N$、训练 token 数 $D$ 与训练计算量 $C$ 放入统一框架，跨多个数量级拟合 loss 曲线\cite{kaplan2020}。其代表性指数约为 $\alpha_N=$\textbf{0.076}、$\alpha_D=$\textbf{0.095}、$\alpha_C=$\textbf{0.050}，并得到早期 compute-optimal 结论：预算增长时模型规模应比数据量增长得更快，近似为
\begin{equation}
  N_{\mathrm{opt}}\propto C^{0.73},\qquad
  D_{\mathrm{opt}}\propto C^{0.27}.
  \label{eq:kaplan-opt}
\end{equation}
这一结论影响了 GPT-3 式“大模型、相对少 token、较早停止”的训练路线\cite{brown2020}。它的贡献不在于给出永恒指数，而在于把“扩大模型是否有效”变成了“可用 pilot 预测目标 run”的工程问题。

\subsection{为什么幂律有用，却不是自然定律}

幂律可以来自不同 regime 下的方差限制、分辨率限制或逼近误差；理论工作能够解释常见的函数形态，但仍不能给出跨数据分布与架构通用的指数\cite{bahri2021}。经验上还会出现 plateau、double descent、迟发跃迁或 slope change，smoothly broken power law 往往比单一幂律更合适\cite{caballero2022}。因此需要把式~\eqref{eq:single-power} 理解为\term{有限区间、固定 recipe 下的局部模型}。

\section{核心数学：从联合损失面到 Compute-Optimal}

\subsection{变量口径与训练计算近似}

表~\ref{tab:symbols} 给出本文统一口径。对于标准 dense decoder-only Transformer，常用一阶近似
\begin{equation}
  C\approx 6ND.
  \label{eq:compute}
\end{equation}
这里 forward 约为 $2ND$ FLOPs，backward 约为 forward 的两倍。该式只适合粗略规划：embedding、长上下文 attention 的 $O(S^2)$ 项、MoE routing、激活专家数和通信都会使其失真。

\begin{table}[htbp]
\centering
\caption{Scaling study 的核心变量与工程口径。}
\label{tab:symbols}
\begin{tabularx}{\textwidth}{C{1.0cm}YY}
\toprule
符号 & 含义 & 实践口径 \\
\midrule
$N$ & 模型规模 & dense 常用 non-embedding parameters；MoE 另记 total/active \\
$D$ & 训练数据量 & 实际送入优化器的 token；重复数据同样计数 \\
$C$ & 训练计算量 & theoretical FLOPs、device FLOPs 或 wall-clock 应分开记录 \\
$L$ & held-out loss & 固定 tokenizer、固定验证分布上的 token-level NLL \\
$E$ & irreducible loss & 需要拟合的渐近项，不应随意固定为论文数字 \\
\bottomrule
\end{tabularx}
\end{table}

\subsection{Chinchilla 联合损失模型及推导}

Hoffmann 等人训练了 400 多个、70M--16B 参数、5B--500B token 的模型，采用联合损失面\cite{hoffmann2022}：
\begin{equation}
  L(N,D)=E+\frac{A}{N^{\alpha}}+\frac{B}{D^{\beta}},
  \qquad A,B,\alpha,\beta>0.
  \label{eq:chinchilla}
\end{equation}
在式~\eqref{eq:compute} 的计算约束下，令 $D=C/(6N)$，代回式~\eqref{eq:chinchilla}：
\begin{equation}
  L(N\mid C)=E+A N^{-\alpha}
  +B\left(\frac{6N}{C}\right)^{\beta}.
  \label{eq:constrained-loss}
\end{equation}
令式~\eqref{eq:constrained-loss} 对 $N$ 的导数为零，可得两个边际收益平衡，并得到
\begin{align}
  N_{\mathrm{opt}}(C)
  &=G\left(\frac{C}{6}\right)^{\frac{\beta}{\alpha+\beta}},
  &G&=\left(\frac{\alpha A}{\beta B}\right)^{\frac{1}{\alpha+\beta}},
  \label{eq:nopt}\\
  D_{\mathrm{opt}}(C)
  &=G^{-1}\left(\frac{C}{6}\right)^{\frac{\alpha}{\alpha+\beta}}.
  \label{eq:dopt}
\end{align}
若记 $a=\beta/(\alpha+\beta)$、$b=\alpha/(\alpha+\beta)$，则 $a+b=1$。直觉上，数据项的边际收益越高，最优点越向数据侧移动；模型项的边际收益越高，最优点越向参数侧移动。

\subsection{Kaplan 与 Chinchilla 的差异}

\begin{table}[htbp]
\centering
\caption{两代 compute-optimal 结论的差异。}
\label{tab:kaplan-chinchilla}
\begin{tabularx}{\textwidth}{Y C{3.2cm} C{3.2cm}}
\toprule
比较项 & Kaplan 2020 & Chinchilla 2022 \\
\midrule
$N_{\mathrm{opt}}$ 对 $C$ 的指数 & 0.73 & 约 0.49--0.50 \\
$D_{\mathrm{opt}}$ 对 $C$ 的指数 & 0.27 & 约 0.50--0.51 \\
工程含义 & 更快增大模型、较早停止 & 参数与 token 近似同比例增长 \\
代表实践 & GPT-3 175B / 300B & Chinchilla 70B / 1.4T \\
\bottomrule
\end{tabularx}
\end{table}

Chinchilla 以与 Gopher 近似的训练计算量，把 280B 参数、300B token 改为 70B 参数、1.4T token，并在其报告的多数任务上更好\cite{rae2021,hoffmann2022}。这使 scaling law 从“解释曲线”变成了“直接改变训练配方”的工具。

“约 \textbf{20 tokens/parameter}”只是特定数据、模型族和实验区间下的工程口径。它优化单次训练 FLOPs，不包含推理成本、显存、数据质量与重复。Besiroglu 等人的复现还发现 Chinchilla Approach 3 的拟合存在优化器提前停止和置信区间过窄问题；重新拟合的点估计仍与约 \textbf{20--26 tokens/parameter} 相容，但\textit{合理区间明显更宽}\cite{besiroglu2024}。稳健结论是\term{参数和 token 应共同扩展}，而不是精确常数可跨体系迁移。

\section{技术演进：Scaling 的自变量不断增加}

\subsection{迁移、多模态与生产级预测}

Henighan 等人把幂律扩展到图像、视频、多模态和数学问题生成\cite{henighan2020}；Hernandez 等人把预训练收益解释为下游数据的“有效倍增”，并刻画 transfer scaling\cite{hernandez2021}。GPT-4 技术报告展示了生产级用途：用计算量低至四个数量级的小模型预测最终内部 codebase loss，并以低至三个数量级的模型预测 HumanEval 子集表现；报告同时明确部分能力仍难预测\cite{openai2023gpt4}。这支持的是\term{受控训练体系内的可预测性}，不是任意能力的普适预测。

\subsection{从平滑 Loss 到“涌现”争论}

“Emergent abilities”把小模型近乎为零、大模型突然上升的 benchmark 能力解释为跃迁\cite{wei2022}。但不连续指标本身就能制造跳变：将 exact match 换为 token-level accuracy 或 Brier score 后，很多曲线会更平滑\cite{schaeffer2023}。因此至少要区分：模型内部能力可能平滑增长；任务成功率可能因阈值呈 S 型；prompt 与采样制造高方差；三个模型点不足以证明 phase transition。

\subsection{从“更多数据”到“有效数据”}

Data-constrained scaling 研究表明，在固定计算预算下适度重复数据可能损失很小，但重复的边际价值最终趋近于零\cite{muennighoff2023}。DeepSeek LLM 报告不同数据版本会改变最优分配指数，说明 scaling coefficient 本身可作为数据与 recipe 的体系内诊断量\cite{deepseekllm2024}。DataComp-LM 用标准化小模型实验展示过滤策略能显著移动性能前沿\cite{datacomplm2024}；data mixing law 则把 mixture proportion 也作为可拟合变量\cite{datamixing2024}。

把这些结果统一起来，可以写成更现实但不可直接一次拟合的关系：
\begin{equation}
 \mathrm{quality}=f\!\left(N,D,Q_{\mathrm{data}},C_{\mathrm{train}},
 C_{\mathrm{test}},\mathrm{architecture},\mathrm{objective}\right).
 \label{eq:modern-scaling}
\end{equation}
这里 $Q_{\mathrm{data}}$ 不是已解决的单一标量，而是 filtering、dedup、domain mix、污染、信息密度与重复结构的集合。

\subsection{从训练最优到全生命周期最优}

若模型会被调用 $R$ 次，决策目标更接近
\begin{equation}
  J=C_{\mathrm{train}}+R\,c_{\mathrm{infer}}(N,C_{\mathrm{test}}),
  \label{eq:lifecycle}
\end{equation}
而不是只最小化训练 FLOPs。Inference-aware scaling 的结果是：\textbf{高调用量下，训练一个更小但看过更多 token 的模型可能拥有更低总成本}\cite{sardana2024}。Test-time scaling 进一步表明，\textit{自适应分配搜索或 verifier compute，在部分问题上可以比单纯扩大参数更有效}\cite{snell2024}。现代 scaling 因而从 $N$--$D$--$C_{\mathrm{train}}$ 三角形扩展到训练、部署与推理时计算的联合优化。

\section{实现细节：怎样做一套可用于发车决策的 Scaling Study}

\subsection{先定义决策，再定义实验}

一个可用研究必须回答可执行问题，例如：给定 $10^{22}$ FLOPs，选 1B/20B tokens 还是 2B/10B tokens；新的 code filter 是否能在等 FLOPs 下移动 OOD code loss frontier；8K repo-level packing 的收益能否从 300M/1B pilot 外推到 7B；或服务量达到多少时 longer-trained 7B 优于 13B。没有决策目标，最终通常只有好看的拟合曲线，没有训练配方。

\subsection{最小实验矩阵与控制变量}

表~\ref{tab:experiment-matrix} 是面向 7B 目标模型的最小可用设计。规模比绝对数更重要：每条 isoFLOP 至少覆盖最优点两侧，最大 pilot 与目标点之间不宜跨越过多数量级。

\begin{table}[htbp]
\centering
\caption{建议的 pilot 实验矩阵。每个格点代表一个完整 run，不是同一 run 的 checkpoint。}
\label{tab:experiment-matrix}
\begin{tabularx}{\textwidth}{Y Y Y Y}
\toprule
维度 & 推荐取值 & 目的 & 最低控制要求 \\
\midrule
模型规模 $N$ & 150M, 300M, 600M, 1.2B, 2.4B & 覆盖约 1 个数量级 & 相同 depth/width 规则与架构族 \\
IsoFLOP $C$ & 4--6 个计算档 & 找到每档 U 型 loss--$N$ 曲线 & 每档至少 4 个 $N/D$ 组合 \\
数据版本 & baseline/candidate & 判断数据策略是否移动 frontier & 同 raw snapshot、tokenizer 与 mix \\
随机性 & 关键点 2--3 seeds & 估计 run variance & 相同初始化分布和数据顺序规则 \\
验证集 & in-domain + OOD + code/repo & 区分拟合与迁移 & 时间切分、去污染、固定版本 \\
\bottomrule
\end{tabularx}
\end{table}

每个 run 至少记录：模型总参数、non-embedding 参数、MoE active 参数、实际 token、unique token、epoch、理论 FLOPs、实测吞吐、序列长度、tokenizer、数据版本、语言/code mix、optimizer、LR schedule、seed、checkpoint loss 与硬件异常。重复 token 和 unique token 必须分开；MoE 的 total/active parameters 必须分开。

\subsection{鲁棒拟合、外推与不确定性}

对观测点 $i$，可以拟合 log-space 残差
\begin{equation}
 r_i=\log\!\left(E+A N_i^{-\alpha}+B D_i^{-\beta}\right)-\log L_i,
 \label{eq:residual}
\end{equation}
并使用 Huber loss。推荐做法是对 $A,B,E$ 使用对数参数化，给 $\alpha,\beta$ 正值边界，多组初始化，保留并标注失败点；bootstrap 必须按 run 重采样，不能把同一 run 的高度相关 checkpoints 当成独立样本。验证集应整条留出更高 compute frontier，不能只随机拆分观测点。

\noindent\textbf{代码 1：Chinchilla 风格联合损失面的最小拟合骨架。}\par\vspace{5pt}
\begin{small}
\setlength{\fboxsep}{7pt}
\setlength{\fboxrule}{0.4pt}
\begin{lrbox}{\fitcodebox}
\begin{minipage}{\dimexpr\textwidth-2\fboxsep-2\fboxrule\relax}
\vspace{1pt}
\begin{verbatim}
import numpy as np
from scipy.optimize import least_squares

def residual(theta, n, d, observed_loss):
    log_a, log_b, log_e, alpha, beta = theta
    pred = (np.exp(log_e)
            + np.exp(log_a) * n ** (-alpha)
            + np.exp(log_b) * d ** (-beta))
    return np.log(pred) - np.log(observed_loss)

starts = [[6.0, 6.0, 0.5, a, b]
          for a in (0.15, 0.30, 0.50, 0.80)
          for b in (0.15, 0.30, 0.50, 0.80)]
fits = [least_squares(
            residual, x0=s, args=(n, d, loss),
            bounds=([-20, -20, -20, 0.01, 0.01],
                    [ 40,  40,   5, 2.00, 2.00]),
            loss="huber", f_scale=1e-3)
        for s in starts]
fit = min(fits, key=lambda x: np.sum(x.fun ** 2))
\end{verbatim}
\vspace{-5pt}
\end{minipage}
\end{lrbox}
\noindent\makebox[\textwidth][l]{\fcolorbox{rulegray}{softgray}{\usebox{\fitcodebox}}}\par
\end{small}

\subsection{三种方法互相校验}

\begin{table}[htbp]
\centering
\caption{Compute-optimal 的三种估计方法。}
\label{tab:fit-methods}
\begin{tabularx}{\textwidth}{Y Y Y Y}
\toprule
方法 & 做法 & 优点 & 主要风险 \\
\midrule
Training-curve minimum & 每个 $N$ 跑足，找给定 compute 的最低 loss & 直观 & 需要密集 checkpoint，早停误差大 \\
IsoFLOP profile & 固定多个 $C$，扫描不同 $N/D$ & 直接回答资源分配 & frontier 点少时插值敏感 \\
Parametric surface & 拟合 $L(N,D)$ 后解析求最优 & 利用全部点并连续外推 & 依赖函数形式、初值和异常点 \\
\bottomrule
\end{tabularx}
\end{table}

\textbf{三种方法在目标尺度附近一致，且 scale holdout 落入预测区间，才适合发车昂贵的大 run。}必做诊断包括：$\log(L-E)$ 对 $\log N/\log D/\log C$；每条 isoFLOP 的 U 型曲线；residual heatmap；预测值对实际值；$N_{\mathrm{opt}}(C)$ 与 $D_{\mathrm{opt}}(C)$ 的 bootstrap band；目标 run 的 prediction interval。

当 residual 随 scale 呈系统趋势、最大尺度持续偏向同一方向、$E$ 贴边、数据或 tokenizer 已改变、dense 切换 MoE、context 增长导致 attention FLOPs 不可忽略时，应判定旧 law 失效并重新拟合。

\section{公开 LLM 尺寸：实践如何印证、又如何混淆 Scaling}

附录~\ref{app:model-inventory} 收录 2020--2026 年 28 个公开模型家族。主文只抽取适合形成判断的对照。表~\ref{tab:dense-families} 中，GPT-3、Pythia、Cerebras-GPT 尺寸点密集；LLaMA、Llama 3.1、Qwen2.5/Qwen3 与 OLMo 2 更能代表现代 data-rich 与 deployment-oriented 训练。

\begin{table}[htbp]
\centering
\small
\caption{具有代表性的 dense 模型家族及其证据边界。}
\label{tab:dense-families}
\begin{tabularx}{\textwidth}{p{2.2cm}p{3.9cm}p{2.5cm}Y}
\toprule
家族 & 尺寸梯度 & 训练 token & 能支持的判断 \\
\midrule
GPT-3 & 125M--175B，共 8 点 & 每个约 300B & 同 recipe 的 size/few-shot 趋势；不适合估 compute-optimal ratio\cite{brown2020} \\
Pythia & 70M--12B，共 8 点 & 每个约 300B & 开放 checkpoints 与固定数据顺序，适合训练动力学\cite{biderman2023} \\
Cerebras-GPT & 111M--13B，共 7 点 & 20 tokens/param. & 公开 Chinchilla-style family\cite{cerebrasgpt2023} \\
LLaMA & 7B, 13B, 33B, 65B & 1.0T/1.4T & data-rich 小模型可超过旧大模型；token budget 不同\cite{llama2023} \\
Llama 3.1 & 8B, 70B, 405B & 超过 15T & 现代家族远超 20 tokens/parameter；不可与旧代连线\cite{llama31} \\
Qwen2.5 & 0.5B--72B，共 7 点 & 18T & 现代 dense 尺寸梯度，说明数据/recipe 移动 frontier\cite{qwen25} \\
Qwen3 Dense & 0.6B--32B，共 6 点 & 约 36T & 新代小模型追平旧代大模型是曲线整体移动，不是参数失效\cite{qwen3} \\
OLMo 2 & 1B, 7B, 13B, 32B & 最高 5T--6T & 数据、代码、recipe 与 checkpoints 可审计\cite{olmo2} \\
\bottomrule
\end{tabularx}
\end{table}

公开实践给出五条相对稳健的结论。第一，\textbf{参数不是独立自变量}；LLaMA-13B 超过 GPT-3 175B 的多项结果来自数据、token、架构和 recipe 共同移动曲线。第二，\textbf{训练显著 data-rich}：从 GPT-3 \textbf{175B/300B}，到 Chinchilla \textbf{70B/1.4T}，再到 Llama 3.1 的 \textbf{15T+} 与 Qwen3 的约 \textbf{36T}。第三，\textit{这不表示 Chinchilla“错误”}，因为现代模型还优化推理成本、可部署性和模型家族复用。第四，\textbf{过滤、混合与领域数据能够横向移动整条 frontier}。第五，\textbf{MoE 把“规模”拆成容量与计算}。

\subsection{MoE：Total Parameters 与 Active Parameters 是两条轴}

表~\ref{tab:moe} 展示如果只看 total parameters，会如何误读 MoE。总参数更接近可存储的专家容量，active parameters 更接近单 token 的前向计算，但 routing、共享专家、通信与负载均衡仍不可忽略。

\begin{table}[htbp]
\centering
\small
\caption{代表性 MoE 模型。Total 与 active parameters 不能混为一谈。}
\label{tab:moe}
\begin{tabularx}{\textwidth}{Y R{2.0cm} R{2.0cm} R{1.7cm} Y}
\toprule
模型 & Total & Active/token & Tokens & 正确解读 \\
\midrule
Mixtral 8x7B & 46.7B & 12.9B & 未公开 & 存储容量近 47B，单 token 计算近 13B\cite{mixtral2024} \\
DeepSeek-V2 & 236B & 21B & 8.1T & 大 total capacity 与低 active compute\cite{deepseekv2} \\
DeepSeek-V3 & 671B & 37B & 14.8T & 不能将 671B 直接代入 dense 式~\eqref{eq:compute}\cite{deepseekv3} \\
Qwen3-30B-A3B & 30B & 3B & 约 36T & 小 active 模型借专家容量提升效果\cite{qwen3} \\
Qwen3-235B-A22B & 235B & 22B & 约 36T & total/active 共同决定容量与成本\cite{qwen3} \\
Llama 4 Scout/Maverick & 109B/400B & 17B/17B & 未公开 & 多模态与 distillation 混杂，不是 clean scaling suite\cite{llama4} \\
DeepSeek-V4 Flash/Pro & 284B/1.6T & 13B/49B & 未公开 & 2026 endpoint，只支持架构趋势观察\cite{deepseekv4} \\
\bottomrule
\end{tabularx}
\end{table}

\section{对 Code Pretraining 与数据 Pipeline 的直接启示}

\subsection{不要只拟合 Aggregate Loss}

代码数据策略至少要分别拟合
\begin{equation}
  L_{\mathrm{general}}(N,D),\quad L_{\mathrm{code}}(N,D),\quad
  L_{\mathrm{repo}}(N,D),\quad L_{\mathrm{swe}}(N,D).
  \label{eq:domain-losses}
\end{equation}
同一数据改动可能降低 code loss，却提高 general loss；aggregate loss 会被占比最大的 web text 掩盖。loss frontier 与 HumanEval、MBPP、CrossCodeEval、SWE held-out 应分开报告，下游指标还要完成污染审计。

\subsection{比较数据策略时必须等价的 Setting}

对 baseline 与 candidate pipeline，应固定 raw snapshot、tokenizer、unique token budget、语言/code mix、$N/D/C$、batch tokens 与 LR schedule；记录过滤保留率、dedup ratio、平均文档长度、code token ratio、repo coverage；同时评估 in-domain 与时间切分 OOD validation。核心指标不只包括最终 $\Delta L$，还应包括每 $10^{20}$ FLOPs 的 loss 改善，以及达到目标 loss 所需 FLOPs。

若 candidate 用 $D$ token 达到 baseline 用 $kD$ token 才达到的 loss，则可定义 data multiplier
\begin{equation}
  M_{\mathrm{data}}=k.
  \label{eq:data-multiplier}
\end{equation}
它能把“loss 降低 0.01”翻译为 token efficiency 与算力收益，但前提是两个版本的 tokenizer、验证集与训练轨迹可比。

\subsection{Repo-Level 组织带来的特殊混杂}

repo/func-level packing 同时改变 token 邻接分布、cross-file mutual information、sequence utilization、样本长度、attention FLOPs、重复函数密度与污染传播范围。因此“训练 token 相同”仍不等价；还要报告 effective non-padding tokens、平均 attention length、每 token 实测 FLOPs，以及 file/repo 级 dedup 后的 unique coverage。

推荐两阶段闭环：\textbf{第一阶段以 150M--2.4B 覆盖 4--6 条 isoFLOP，筛选 pipeline 并预测 7B}；\textbf{第二阶段至少做一个 7B 等级、等训练 token 的 baseline/candidate A/B run}。发车门槛可设为：目标尺度 predicted code OOD loss 改善的 \textbf{80\% bootstrap lower bound 大于 0}；general OOD loss 不越 guardrail；等 FLOPs 的 token efficiency 可量化；语言配比与污染无显著偏移；target confirmation 落在 pilot prediction interval 内。

\section{判断框架：Scaling Law 能预测什么、不能预测什么}

\begin{table}[htbp]
\centering
\caption{不同预测目标的相对可信度。}
\label{tab:predictability}
\begin{tabularx}{\textwidth}{Y C{2.0cm} Y}
\toprule
目标 & 可信度 & 主要原因 \\
\midrule
同 tokenizer、同分布 validation loss & 高 & 连续、低噪声、与训练目标一致 \\
邻近尺度 compute-optimal $N/D$ & 中高 & 需要足够 frontier 点和误差带 \\
OOD domain loss & 中 & 取决于 domain overlap 与 data mix \\
连续 downstream metric & 中低 & prompt、transfer 与 contamination 介入 \\
Exact match/pass@1 & 低 & 离散阈值与高方差 \\
新能力何时出现 & 很低 & 定义、指标、数据和 elicitation 均不稳定 \\
\bottomrule
\end{tabularx}
\end{table}

审阅任何 scaling claim 时，应先回答三个一级问题：\textbf{横轴究竟是什么？纵轴究竟是什么？控制变量是否固定？}再检查覆盖多少数量级；$E$ 如何拟合；是否有最大尺度 holdout；是否报告 seed variance、bootstrap interval 与 residual；是否区分 unique/repeated tokens；污染、prompt 与评分器版本是否受控；优化目标究竟是训练 FLOPs、wall-clock、推理成本还是总拥有成本。

常见错误包括：把模型排行榜当 scaling law；把 20 tokens/parameter 当硬规范；比较不同 tokenizer 的 nats/token；随机拆 checkpoints 做验证；只报 $R^2$；以 total parameters 估 MoE FLOPs；把 benchmark 跳变直接解释为机制突变；静默删除失败 run；把论文系数直接搬到内部数据。\textbf{通常可迁移的是函数形式与实验方法，不是系数。}

\section{总结与展望}

LLM scaling law 的技术史可以概括为三次认知升级：Hestness 与 Kaplan 说明性能随资源增长具有可预测结构；Chinchilla 把问题从“是否变好”改为“参数与数据如何分配”；此后数据质量、重复、mixture、MoE active compute、推理时计算与部署经济性成为新的扩展轴。

对入门者，最重要的认知是：\textbf{scaling law 不是智能增长的自然定律，而是固定体系内的局部经验模型。}对预训练团队，最重要的动作是：\textbf{在自己的 tokenizer、数据版本、模型架构和 OOD validation 上做 pilot}，用 isoFLOP 与 parametric surface 互相校验，报告外推区间，再用目标尺度 A/B run 闭环。\textit{公开模型实践提供方向性印证，但不能替代内部 scaling study。}

仍待解决的问题包括：如何把数据质量压缩为可跨体系比较的变量；多域 mixture 如何随规模联合外推；合成数据的有效 token 如何计量；MoE 的容量、active compute、routing entropy 与通信如何统一；train-time 与 test-time compute 如何联合最优；token loss 如何映射到 repo-level editing 和 agent 长程任务；以及何时能在代价可控时提前检测 regime change。下一阶段更值得研究的不是一条更宏大的“万能曲线”，而是\term{带证据等级、误差边界与决策闭环的局部 scaling laws}。

\appendix
\clearpage
\section{术语速查}
\label{app:glossary}

\begin{table}[htbp]
\centering
\caption{Scaling law 常用术语。}\label{tab:glossary}
\begin{tabularx}{\textwidth}{p{3.6cm}X}
\toprule
术语 & 定义 \\
\midrule
Scaling law & 资源规模与 loss/性能之间的经验函数关系。 \\
Power law & $y=A x^{-\alpha}$；减去渐近项后在 log--log 坐标近似直线。 \\
Reducible loss & 可通过更多模型、数据或计算降低的 loss。 \\
Irreducible loss & 当前数据分布与模型族下的渐近下限。 \\
Compute-optimal & 固定训练计算预算时使目标 loss 最低的资源分配。 \\
IsoFLOP & 固定训练 FLOPs，扫描不同模型/数据组合。 \\
Tokens/parameter & 训练 token 数与参数量之比；不是跨 recipe 通用常数。 \\
Active parameters & MoE 中每个 token 实际激活并参与计算的参数。 \\
Data multiplier & candidate 数据达到同 loss 时相当于 baseline 的有效数据倍数。 \\
Broken scaling law & 允许不同 scale regime 中斜率平滑变化的经验规律。 \\
\bottomrule
\end{tabularx}
\end{table}

\section{公开模型尺寸与 Scaling 证据清单}
\label{app:model-inventory}

表~\ref{tab:model-inventory} 的 token 数来自对应论文或官方发布页。“未公开”表示没有 fit-ready 的公开数字，并非推断为零；open weights 不等同于训练数据与过程完全开放。若 token budget 随尺寸变化，应回到原报告逐模型展开。2026 年模型只在官方资料明确给出尺寸时纳入，并统一标为 C 级趋势证据。

\scriptsize
\hbadness=10000
\setlength{\tabcolsep}{1.7pt}
\setlength{\LTleft}{0pt}
\setlength{\LTright}{0pt}
\noindent\begin{minipage}{\textwidth}
\captionof{table}{2020--2026 年公开 LLM 家族、尺寸、训练 token 与 scaling 证据。}\label{tab:model-inventory}
\end{minipage}
\begin{longtable}{C{0.55cm}p{1.50cm}p{1.10cm}p{3.20cm}p{2.00cm}p{1.80cm}C{0.65cm}p{3.80cm}}
\toprule
年份 & 家族 & 架构 & Total parameter sizes & Active/token & Pretraining tokens & 证据 & 可支持的判断 \\
\midrule
\endfirsthead
\multicolumn{8}{l}{\small 表~\thetable（续）}\\
\toprule
年份 & 家族 & 架构 & Total parameter sizes & Active/token & Pretraining tokens & 证据 & 可支持的判断 \\
\midrule
\endhead
\midrule
\multicolumn{8}{r}{\small 下页续}\\
\endfoot
\bottomrule
\endlastfoot
2020 & GPT-3\cite{brown2020} & Dense & 125M; 350M; 760M; 1.3B; 2.7B; 6.7B; 13B; 175B & 同 total & 每个约 300B & A & 同家族参数 scaling 与 few-shot 趋势；不能估 compute-optimal ratio。 \\
2021 & Gopher\cite{rae2021} & Dense & 44M; 117M; 417M; 1.4B; 7.1B; 280B & 同 total & 主模型约 300B & A$-$ & 尺寸 scaling 与不同任务收益不均衡。 \\
2022 & Chinchilla\cite{hoffmann2022} & Dense & 70B & 同 total & 1.4T & A & 与 Gopher 的同计算对照及 compute-optimal allocation。 \\
2022 & PaLM\cite{palm2022} & Dense & 8B; 62B; 540B & 同 total & 780B & A$-$ & 家族内尺寸趋势与 apparent benchmark thresholds。 \\
2022 & OPT\cite{opt2022} & Dense & 125M; 350M; 1.3B; 2.7B; 6.7B; 13B; 30B; 66B; 175B & 同 total & 主设置约 180B & B$+$ & 尺寸覆盖广；须结合 recipe 与 loss 控制。 \\
2022 & BLOOM\cite{bloom2022} & Dense & 560M; 1.1B; 1.7B; 3B; 7.1B; 176B & 同 total & 176B 模型约 350B & B & 多语言规模观察；尺寸梯度有缺口。 \\
2023 & Pythia\cite{biderman2023} & Dense & 70M; 160M; 410M; 1B; 1.4B; 2.8B; 6.9B; 12B & 同 total & 每个约 300B & A$+$ & 最佳公开受控 suite 之一，适合训练动力学。 \\
2023 & Cerebras-GPT\cite{cerebrasgpt2023} & Dense & 111M; 256M; 590M; 1.3B; 2.7B; 6.7B; 13B & 同 total & 20 tokens/parameter & A & 公开 Chinchilla-style family。 \\
2023 & LLaMA\cite{llama2023} & Dense & 7B; 13B; 33B; 65B & 同 total & 1.0T 或 1.4T & B$+$ & data-rich 小模型可超过旧大模型；不是单轴拟合。 \\
2023 & Llama 2\cite{llama2} & Dense & 7B; 13B; 70B & 同 total & 2T & B$+$ & 同代部署梯度；点数不足以稳健拟合指数。 \\
2023 & Falcon\cite{falcon2023} & Dense & 7B; 40B; 180B & 同 total & 180B 最高约 3.5T & B & data-rich 与 web filtering；token budget 随尺寸变化。 \\
2023 & DeepSeek LLM\cite{deepseekllm2024} & Dense & 7B; 67B & 同 total & 2T & A$-$ & pilot 外推，并显示数据版本会改变拟合指数。 \\
2023 & Mixtral 8x7B\cite{mixtral2024} & MoE & 46.7B & 12.9B & 未公开 & C$+$ & total 与 active compute 必须分开；不是 scaling fit。 \\
2024 & Open\allowbreak ELM\cite{openelm2024} & Dense & 270M; 450M; 1.1B; 3B & 同 total & 1.5T & B$+$ & 同 token budget 的小模型架构比较。 \\
2024 & OLMo\cite{olmo2024} & Dense & 1B; 7B & 同 total & 约 2T--3T & B$+$ & 数据顺序、checkpoints 与训练日志可审计。 \\
2024 & Llama 3.1\cite{llama31} & Dense & 8B; 70B; 405B & 同 total & 超过 15T & B$+$ & 现代 data-rich family；不支持 universal coefficients。 \\
2024 & Qwen2.5\cite{qwen25} & Dense & 0.5B; 1.5B; 3B; 7B; 14B; 32B; 72B & 同 total & 18T & B$+$ & 现代尺寸梯度，显示 data/recipe 移动 frontier。 \\
2024 & Qwen2.5-Coder\cite{qwen25coder} & Dense & 0.5B; 1.5B; 3B; 7B; 14B; 32B & 同 total & 续训语料 5.5T & B$+$ & 代码域梯度；需区分 base pretraining 与 CPT。 \\
2024 & Gemma 1\cite{gemma1} & Dense & 2B; 7B & 同 total & 2T; 6T & B & 两个点；说明部署家族不固定 tokens/parameter。 \\
2024 & Gemma 2\cite{gemma2} & Dense & 2B; 9B; 27B & 同 total & 2T; 8T; 13T & C$+$ & 2B/9B 含 distillation，反例说明不能只看参数。 \\
2024 & DeepSeek-V2\cite{deepseekv2} & MoE & 236B & 21B & 8.1T & B & 分离存储容量与 per-token compute。 \\
2024 & DeepSeek-V3\cite{deepseekv3} & MoE & 671B & 37B & 14.8T & B & 大规模 MoE 与计算披露；单 endpoint 不是 law。 \\
2024 & OLMo 2\cite{olmo2} & Dense & 1B; 7B; 13B; 32B & 同 total & 最高 5T--6T & B$+$ & 现代 fully open family，适合数据质量研究。 \\
2025 & Qwen3\cite{qwen3} & Dense+\allowbreak MoE & 0.6B; 1.7B; 4B; 8B; 14B; 32B; 30B-MoE; 235B-MoE & Dense 同 total; MoE 3B/22B & 约 36T & B$+$ & 同代 dense 梯度与 active-parameter MoE 对照。 \\
2025 & Llama 4\cite{llama4} & MoE & Scout 109B; Maverick 400B & 均 17B & 未以 fit-ready 形式披露 & C$+$ & 多模态与 distillation 混杂，只支持 active/total 区分。 \\
2025 & Gemma 3\cite{gemma3} & Dense 多模态 & 270M; 1B; 4B; 12B; 27B & 同 total & 随尺寸变化 & C$+$ & 部署尺寸梯度；多模态与 recipe 限制经典比较。 \\
2026 & Gemma 4\cite{gemma4} & Dense+\allowbreak MoE 多模态 & E2B; E4B; 12B; 31B; 26B-A4B & A4B 约 4B active & 未公开 & C & 当前 inventory；缺乏拟合 scaling law 的证据。 \\
2026 & DeepSeek-V4\cite{deepseekv4} & MoE & Flash 284B; Pro 1.6T & 13B; 49B & 未公开 & C & total capacity 增长而 active compute 相对克制；非 fit-ready。 \\
\end{longtable}

\clearpage
\section{建议的数据记录 Schema 与发车清单}
\label{app:schema}

\begin{table}[htbp]
\centering
\caption{Scaling experiment 最低记录字段。}\label{tab:schema}
\begin{tabularx}{\textwidth}{p{4.0cm}X}
\toprule
字段组 & 必填内容 \\
\midrule
模型 & total params、non-embedding params、active params、layers、hidden size、heads、experts、context length \\
数据 & raw snapshot、pipeline version、tokenizer hash、actual/unique tokens、epochs、语言与 code mix、dedup ratio \\
训练 & global batch tokens、optimizer、peak/min LR、warmup、schedule、weight decay、seed、precision、gradient clipping \\
计算 & theoretical FLOPs、measured FLOPs、tokens/s、GPU/NPU hours、MFU、失败与重启记录 \\
评估 & validation dataset hash、in/OOD loss、bits/byte、benchmark harness、prompt、sampling、污染审计版本 \\
拟合 & 函数形式、初值、bounds、loss function、异常点策略、run-level bootstrap、holdout frontier、prediction interval \\
决策 & 目标规模、预算、$N/D$ 推荐区间、guardrail、预期 data multiplier、confirmation run 与 go/no-go 结论 \\
\bottomrule
\end{tabularx}
\end{table}

\noindent\textbf{发车前检查：}
\begin{enumerate}
  \item 三种 compute-optimal 方法是否在目标尺度附近一致？
  \item 最大 compute holdout 是否落入 prediction interval？
  \item 目标决策对 $E,\alpha,\beta$ 的不确定性是否稳健？
  \item 数据版本、tokenizer、架构或 context 是否发生 regime change？
  \item OOD/code/repo loss 与 aggregate loss 是否给出一致或可解释的信号？
  \item 推理量、显存和 latency 是否使 training-optimal 偏离 lifecycle-optimal？
\end{enumerate}

\clearpage
\phantomsection
\addcontentsline{toc}{section}{参考文献}
\begin{thebibliography}{99}
\small
\setlength{\itemsep}{1pt}
\setlength{\parskip}{0pt}
\hbadness=10000
\newcommand{\bibentry}[5]{#1. \href{#5}{\textit{#2}}. #3, #4.}

\bibitem{hestness2017} \bibentry{J. Hestness et al.}{Deep Learning Scaling is Predictable, Empirically}{arXiv:1712.00409}{2017}{https://arxiv.org/abs/1712.00409}
\bibitem{rosenfeld2019} \bibentry{J. S. Rosenfeld et al.}{A Constructive Prediction of the Generalization Error Across Scales}{ICLR}{2020}{https://arxiv.org/abs/1909.12673}
\bibitem{kaplan2020} \bibentry{J. Kaplan et al.}{Scaling Laws for Neural Language Models}{arXiv:2001.08361}{2020}{https://arxiv.org/abs/2001.08361}
\bibitem{bahri2021} \bibentry{Y. Bahri et al.}{Explaining Neural Scaling Laws}{JMLR}{2024}{https://arxiv.org/abs/2102.06701}
\bibitem{caballero2022} \bibentry{E. Caballero et al.}{Broken Neural Scaling Laws}{arXiv:2210.14891}{2022}{https://arxiv.org/abs/2210.14891}
\bibitem{brown2020} \bibentry{T. B. Brown et al.}{Language Models are Few-Shot Learners}{NeurIPS}{2020}{https://arxiv.org/abs/2005.14165}
\bibitem{hoffmann2022} \bibentry{J. Hoffmann et al.}{Training Compute-Optimal Large Language Models}{NeurIPS}{2022}{https://doi.org/10.48550/arXiv.2203.15556}
\bibitem{rae2021} \bibentry{J. W. Rae et al.}{Scaling Language Models: Methods, Analysis \& Insights from Training Gopher}{arXiv:2112.11446}{2021}{https://arxiv.org/abs/2112.11446}
\bibitem{besiroglu2024} \bibentry{T. Besiroglu et al.}{Chinchilla Scaling: A Replication Attempt}{arXiv:2404.10102}{2024}{https://arxiv.org/abs/2404.10102}
\bibitem{henighan2020} \bibentry{T. Henighan et al.}{Scaling Laws for Autoregressive Generative Modeling}{arXiv:2010.14701}{2020}{https://arxiv.org/abs/2010.14701}
\bibitem{hernandez2021} \bibentry{D. Hernandez et al.}{Scaling Laws for Transfer}{arXiv:2102.01293}{2021}{https://arxiv.org/abs/2102.01293}
\bibitem{openai2023gpt4} \bibentry{OpenAI}{GPT-4 Technical Report}{arXiv:2303.08774}{2023}{https://arxiv.org/abs/2303.08774}
\bibitem{wei2022} \bibentry{J. Wei et al.}{Emergent Abilities of Large Language Models}{TMLR}{2022}{https://arxiv.org/abs/2206.07682}
\bibitem{schaeffer2023} \bibentry{R. Schaeffer, B. Miranda, and S. Koyejo}{Are Emergent Abilities of Large Language Models a Mirage?}{NeurIPS}{2023}{https://arxiv.org/abs/2304.15004}
\bibitem{muennighoff2023} \bibentry{N. Muennighoff et al.}{Scaling Data-Constrained Language Models}{arXiv:2305.16264}{2023}{https://arxiv.org/abs/2305.16264}
\bibitem{deepseekllm2024} \bibentry{DeepSeek-AI et al.}{DeepSeek LLM: Scaling Open-Source Language Models with Longtermism}{arXiv:2401.02954}{2024}{https://arxiv.org/abs/2401.02954}
\bibitem{datacomplm2024} \bibentry{J. Li et al.}{DataComp-LM: In Search of the Next Generation of Training Sets for Language Models}{arXiv:2406.11794}{2024}{https://arxiv.org/abs/2406.11794}
\bibitem{datamixing2024} \bibentry{J. Ye et al.}{Data Mixing Laws: Optimizing Data Mixtures by Predicting Language Modeling Performance}{arXiv:2403.16952}{2024}{https://arxiv.org/abs/2403.16952}
\bibitem{sardana2024} \bibentry{N. Sardana et al.}{Beyond Chinchilla-Optimal: Accounting for Inference in Language Model Scaling Laws}{ICML}{2024}{https://arxiv.org/abs/2401.00448}
\bibitem{snell2024} \bibentry{C. Snell et al.}{Scaling LLM Test-Time Compute Optimally Can Be More Effective than Scaling Model Parameters}{arXiv:2408.03314}{2024}{https://arxiv.org/abs/2408.03314}
\bibitem{biderman2023} \bibentry{S. Biderman et al.}{Pythia: A Suite for Analyzing Large Language Models Across Training and Scaling}{ICML}{2023}{https://arxiv.org/abs/2304.01373}
\bibitem{cerebrasgpt2023} \bibentry{N. Dey et al.}{Cerebras-GPT: Open Compute-Optimal Language Models Trained on the Pile}{arXiv:2304.03208}{2023}{https://arxiv.org/abs/2304.03208}
\bibitem{llama2023} \bibentry{H. Touvron et al.}{LLaMA: Open and Efficient Foundation Language Models}{arXiv:2302.13971}{2023}{https://arxiv.org/abs/2302.13971}
\bibitem{llama2} \bibentry{H. Touvron et al.}{Llama 2: Open Foundation and Fine-Tuned Chat Models}{arXiv:2307.09288}{2023}{https://arxiv.org/abs/2307.09288}
\bibitem{palm2022} \bibentry{A. Chowdhery et al.}{PaLM: Scaling Language Modeling with Pathways}{arXiv:2204.02311}{2022}{https://arxiv.org/abs/2204.02311}
\bibitem{opt2022} \bibentry{S. Zhang et al.}{OPT: Open Pre-trained Transformer Language Models}{arXiv:2205.01068}{2022}{https://arxiv.org/abs/2205.01068}
\bibitem{bloom2022} \bibentry{T. L. Scao et al.}{BLOOM: A 176B-Parameter Open-Access Multilingual Language Model}{arXiv:2211.05100}{2022}{https://arxiv.org/abs/2211.05100}
\bibitem{falcon2023} \bibentry{E. Almazrouei et al.}{The Falcon Series of Open Language Models}{arXiv:2311.16867}{2023}{https://arxiv.org/abs/2311.16867}
\bibitem{mixtral2024} \bibentry{A. Q. Jiang et al.}{Mixtral of Experts}{arXiv:2401.04088}{2024}{https://arxiv.org/abs/2401.04088}
\bibitem{openelm2024} \bibentry{S. Mehta et al.}{OpenELM: An Efficient Language Model Family with Open Training and Inference Framework}{arXiv:2404.14619}{2024}{https://arxiv.org/abs/2404.14619}
\bibitem{olmo2024} \bibentry{D. Groeneveld et al.}{OLMo: Accelerating the Science of Language Models}{arXiv:2402.00838}{2024}{https://arxiv.org/abs/2402.00838}
\bibitem{llama31} \bibentry{Meta AI}{Introducing Llama 3.1: Our Most Capable Models to Date}{Official release}{2024}{https://ai.meta.com/blog/meta-llama-3-1/}
\bibitem{qwen25} \bibentry{Qwen Team}{Qwen2.5 Technical Report}{arXiv:2412.15115}{2024}{https://arxiv.org/abs/2412.15115}
\bibitem{qwen25coder} \bibentry{B. Hui et al.}{Qwen2.5-Coder Technical Report}{arXiv:2409.12186}{2024}{https://arxiv.org/abs/2409.12186}
\bibitem{gemma1} \bibentry{Gemma Team}{Gemma: Open Models Based on Gemini Research and Technology}{arXiv:2403.08295}{2024}{https://arxiv.org/abs/2403.08295}
\bibitem{gemma2} \bibentry{Gemma Team}{Gemma 2: Improving Open Language Models at a Practical Size}{arXiv:2408.00118}{2024}{https://arxiv.org/abs/2408.00118}
\bibitem{deepseekv2} \bibentry{DeepSeek-AI et al.}{DeepSeek-V2: A Strong, Economical, and Efficient Mixture-of-Experts Language Model}{arXiv:2405.04434}{2024}{https://arxiv.org/abs/2405.04434}
\bibitem{deepseekv3} \bibentry{DeepSeek-AI et al.}{DeepSeek-V3 Technical Report}{arXiv:2412.19437}{2024}{https://arxiv.org/abs/2412.19437}
\bibitem{olmo2} \bibentry{Allen Institute for AI}{OLMo 2: The Best Fully Open Language Model to Date}{Official release}{2024}{https://allenai.org/olmo2}
\bibitem{qwen3} \bibentry{Qwen Team}{Qwen3: Think Deeper, Act Faster}{Official technical blog}{2025}{https://qwenlm.github.io/blog/qwen3/}
\bibitem{llama4} \bibentry{Meta AI}{The Llama 4 Herd: The Beginning of a New Era of Natively Multimodal AI Innovation}{Official release}{2025}{https://ai.meta.com/blog/llama-4-multimodal-intelligence/}
\bibitem{gemma3} \bibentry{Google}{Gemma Models Documentation}{Official documentation; accessed 2026-08-07}{2025}{https://ai.google.dev/gemma/docs/get_started}
\bibitem{gemma4} \bibentry{Google}{Gemma Models Documentation: Current Model Inventory}{Official documentation; accessed 2026-08-07}{2026}{https://ai.google.dev/gemma/docs/get_started}
\bibitem{deepseekv4} \bibentry{DeepSeek}{DeepSeek-V4 Official Release}{Official release}{2026}{https://api-docs.deepseek.com/news/news260424/}

\end{thebibliography}

\end{document}
